% L'espace de noms du DNS : un arbre, découpé en zones déléguées.
\documentclass[border=6pt]{standalone}
\usepackage{sd}
\begin{document}
\begin{tikzpicture}[papier,
    n/.style={bloc=#1, minimum width=24mm, minimum height=7mm},
    level distance=17mm]
  \node[n=Plavande] (racine) at (0,0) {. (racine)};
  \node[n=Pbleu] (com) at (-45mm,-17mm) {com};
  \node[n=Pbleu] (fr) at (0,-17mm) {fr};
  \node[n=Pbleu] (org) at (45mm,-17mm) {org};
  \node[n=Pvert] (ex) at (-45mm,-34mm) {example.com};
  \node[n=Pvert] (wp) at (45mm,-34mm) {wikipedia.org};
  \node[n=Pjaune] (www) at (-62mm,-51mm) {www.example.com};
  \node[n=Pjaune] (api) at (-28mm,-51mm) {api.example.com};
  \node[n=Pjaune] (wfr) at (45mm,-51mm) {fr.wikipedia.org};
  \foreach \a/\b in {racine/com, racine/fr, racine/org, com/ex, org/wp, ex/www, ex/api, wp/wfr}
    \draw[ligne] (\a.south) -- (\b.north);
  % zones déléguées
  \begin{scope}[on background layer]
    \node[cadred, fit=(racine), inner sep=2.5mm] (zr) {};
    \node[cadred, fit=(com), inner sep=2.5mm] (zc) {};
    \node[cadred, fit=(ex)(www)(api), inner sep=2.5mm] (ze) {};
  \end{scope}
  \node[et=Pbrun, anchor=west] at (zr.east) {zone racine : 13 identités de serveurs};
  \node[et=Pbrun, anchor=east] at (zc.west) {zone com};
  \node[et=Pbrun, anchor=north] at (ze.south) {zone example.com : un seul jeu de serveurs};
  \node[eti, anchor=west] at ($(com.south)!0.5!(ex.north)$) {\hspace{2mm}délégation (enregistrements NS)};
\end{tikzpicture}
\end{document}
